\documentclass[10pt,a4paper]{amsart}
\usepackage[margin=19mm]{geometry}
\usepackage{amsmath,amssymb,mathtools}
\usepackage[scr=rsfs,cal=euler]{mathalfa}
\usepackage{xfrac}
\usepackage[unicode,hidelinks,bookmarks=false]{hyperref}
\newcommand{\ie}{i.e.\ }
\newcommand{\cf}{cf.\ }
\newcommand{\resp}{respectively\ }
\newcommand{\Subsection}{\subsection}
\providecommand{\nobreakdash}{\nobreak\hbox{-}}
\setlength{\emergencystretch}{2em}
\multlinegap0pt
\usepackage{amssymb}
\newtheorem{theorem}{Theorem}
\newtheorem{lemma}{Lemma}
\newtheorem{definition}{Definition}
\renewcommand{\geq}{\geqslant}  %package{amssymb}
\renewcommand{\leq}{\leqslant}
\title{\bf Uniform Large Deviations of Mckean-Vlasov Stochastic Fractional $(\alpha,p)$-Laplacian Equations Driven by Superlinear Noise on $\mathbb{R}^d$}
\author{Renhai Wang, Zhang Chen, Bixiang Wang}
\date{Source: arXiv:2607.21862v1 (CC BY 4.0)}
\begin{document}
\maketitle

\noindent\emph{Source and licence.} \bf Uniform Large Deviations of Mckean-Vlasov Stochastic Fractional $(\alpha,p)$-Laplacian Equations Driven by Superlinear Noise on $\mathbb{R}^d$. Version arXiv:2607.21862v1 (CC BY 4.0), distributed by arXiv under the Creative Commons Attribution 4.0 International licence, http://creativecommons.org/licenses/by/4.0/, as recorded in the arXiv licence metadata for this exact version and shown on the abstract page https://arxiv.org/abs/2607.21862v1. Full licence notice: \texttt{licenses/arxiv-2607.21862-CC-BY-4.0.txt}.\par\smallskip

\noindent\textit{Editorial scope.} Counted unit: the article's lemma SCE3 (source0.tex lines 3953--3968). The article's complete original proof of it is delivered with it (source0.tex lines 3969--4276, one \texttt{proof} environment of 308 lines). No mathematics is written, repaired or re-derived: the statement and the proof are the article's own lines, in the article's own notation, and no retained line is substituted or re-flowed.  Retained uncounted as the source-local context the proof needs: lines 186--197; lines 356--370; lines 944--993; lines 1248--1288; lines 1290--1342; lines 1387--1395; lines 1402--1412; lines 1422--1434; lines 1442--1462; lines 2548--2558; lines 2563--2572; lines 3877--3894; lines 3914--3925; lines 3932--3946.  Nothing was omitted from any retained span, so the package carries no omission record.  No retained line contains a citation, so the wrapper carries no bibliography and no bibliography entry.  No retained line contains a figure environment, an image-inclusion command or drawing code, so no image is delivered and the package declares no asset dependency.  Editorial formatting only, in addition: an AMS \texttt{amsart} preamble that restates the article's own declarations and macros used by the retained lines, namely the environments \texttt{theorem}, \texttt{lemma}, \texttt{definition} on separate counters, together with the standard packages those lines need.  The retained environments print in this document as Theorem 1, Lemma 1, Definition 1 in order of appearance, whereas the article prints the same items with the numbering of its own full document.  The complete native source of the article as distributed in the arXiv e-print for this exact version is bundled unchanged as \texttt{sources/205896/source0.tex}.
\begin{equation}\label{Int19-}
\left\{
  \begin{aligned}
  &d \psi(t)+ \mathfrak{L}_{\mathcal{K}_{p}^\alpha} \psi(t) dt
=\sum_{i=1}^2H_i(t,x,\psi(t),\mathcal{L}_{
\psi(t)})dt
+ \sqrt{\epsilon}\sum_{k\in\mathbb{N}} \sigma_{k}(t,x,\psi(t),\mathcal{L}_{
\psi(t)}) d\mathcal{W}_k, \ \ t>0, \\
  &\psi(0,x)=\psi_0(x),\\
  \end{aligned}
\right.
\end{equation}

\begin{theorem}\label{Theorem4}(Global-in-time well-posedness)
Let conditions {\bf A} and {\bf B}
in Section 2 with $p^*\in[2,p)$ and $q^*\in[2,q)$ be satisfied.
If $\psi_0\in L^{2}(\Omega,\mathfrak{F}_0,\mathbb{H})$, then problem \eqref{Int19-} has a unique solution $\psi$ in the sense of Definition \ref{Defsolutionsolution} such that
\begin{align}\label{G8}
&
\mathbf{E}\bigg[\sup_{r\in[0,T]}\|\psi
(r)
\|^{2}
+\int_{0}^{T}\Big(\|\psi(r)
\|_{\mathbb{V}_1}^{p} +\|\psi(r)\|^q_{\mathbb{V}_2}\Big)
 dr\bigg]\leq M(T) (1+\mathbf{E}[\|\psi_0\|^2]),\ \ \forall\ T>0,
\end{align}
where $M(T)>0$ is a constant independent of $\epsilon$.
\end{theorem}

\begin{lemma}(\texttt{Strong monotonicity})
\label{Monotonicity}
Let condition \textbf{A} be satisfied. Then we have:

(i)The operator $\mathbf{A}_1(t,\cdot,\cdot): \mathbb{V}_1\times\mathcal{P}_2
(\mathbb{H})\rightarrow
\mathbb{V}^*_1$ is strongly monotone, that is, for all $t\in\mathbb{R}^+$, $\psi_1,\psi_2\in \mathbb{V}_1$ and $\mu_1,\mu_2\in\mathcal{P}_2
(\mathbb{H})$,
\begin{align*}
& _{\mathbb{V}^*_1}
\big\langle \mathbf{A}_1(t,\psi_1,\mu_1)-
\mathbf{A}_1(t,\psi_2,\mu_2),
\psi_1-\psi_2\big\rangle_{\mathbb{V}_1}
\leq -\frac{1}{2^{p+1}K}\|\psi_1-\psi_2
\|^p_{\dot{\mathbb{V}}_1}
 -2^{-p}\beta\|\psi_1-\psi_2\|^p_p
\nonumber\\& -\frac{1}{8K} \int_{\mathbb{R}^d}\int_{\mathbb{R}^d}
\big( |\psi_1(x)-\psi_1(y)|^{p-2}
+|\psi_2(x)-\psi_2(y)|^{p-2}  \big)
|\psi_1(x)-\psi_1(y)
-(\psi_2(x)-\psi_2(y))|^{2}|x-y|^{-(d+\alpha p)}dxdy
\\&\ \ -\frac{\beta}{4} \int_{\mathbb{R}^d}  \big(|\psi_1(x)|^{p-2}+|\psi_2(x)|^{p-2}\big)
|\psi_1(x)-\psi_2(x)|^2dx
+\|\phi_6(t)\|_{\infty}
\|\psi_1-\psi_2\|^2+
\|\phi_7(t)
\|_{1}
d_{\mathcal{P}_2}^2
(\mu_1,\mu_2).
\end{align*}

(ii)The operator $\mathbf{A}_2(t,\cdot,\cdot): \mathbb{V}_2\times\mathcal{P}_2
(\mathbb{H})\rightarrow
\mathbb{V}^*_2$
is strongly monotone, that is, for all   $t\in\mathbb{R}^+$, $\psi_1,\psi_2\in \mathbb{V}_2$ and $\mu_1,\mu_2\in\mathcal{P}_2
(\mathbb{H})$,
\begin{align*}&
_{\mathbb{V}^*_2}\langle \mathbf{A}_2(t,\psi_1,\mu_1)
-\mathbf{A}_2(t,\psi_2,\mu_2) ,\psi_1-\psi_2\rangle_{\mathbb{V}_2}
 \leq-2^{-q}\widehat{\beta}
\|\psi_1-\psi_2\|_{\mathbb{V}_2}^{q} \\& -\frac{\widehat{\beta}}{4}\int_{\mathbb{R}^d}  \big(|\psi_1(x)|^{q-2}+|\psi_2(x)|^{q-2}\big)
|\psi_1(x)-\psi_2(x)|^2dx
  +\|\widehat{\phi}_6(t)\|_{\infty}
\|\psi_1-\psi_2\|^2+
\|\widehat{\phi}_7(t)
\|_{1}
d_{\mathcal{P}_2}^2
(\mu_1,\mu_2).
\end{align*}
\end{lemma}

\begin{align}\label{h2-}
&M\|\sigma(t,\cdot,\psi,\mu)\|^2_{
\ell^2(\mathbb{N},
\mathbb{H})}
=M\sum_{k\in\mathbb{N}}\int_{\mathbb{R}^d}
|\sigma_k(t,x,\psi(x),\mu)|^2dx
\nonumber\\
&\leq
M\sum_{k\in\mathbb{N}}\int_{\mathbb{R}^d}
\Big(
\sigma_{1,k}(x)
|\psi(x)|^{p^*}+
\sigma_{2,k}(x)
|\psi(x)|^{q^*}
 +
\sigma_{3,k}(x)
\big(1+\mu\big(\|\cdot\|^2\big)\big)\Big)dx
\nonumber\\
&\leq
 M
 \Bigg(
 \sum_{k\in\mathbb{N}}\|\sigma_{1,k} \|
 _{\frac p{p-p^*}}
 \| \psi\|^{p^*}_{p}
 +
 \sum_{k\in\mathbb{N}}\|\sigma_{2,k}\|
 _{\frac q{q-q^*} }
  \| \psi\|^{q^*}_{q}
 +
 \sum_{k\in\mathbb{N}}\|\sigma_{3,k}\|
 _{1}\big(1+\mu\big(\|\cdot\|^2
\big)\big)\Bigg),
\nonumber\\
&\le
\varepsilon_1
\|\psi\|_p^{p}+\varepsilon_2
\|\psi\|_{\mathbb{V}_2}^{q}\
+C_2(M,\varepsilon_1,\varepsilon_2)
\big(1+\mu\big(\|\cdot\|^2
\big)\big),
\end{align}

\begin{align}\label{h2--}
&M\|\sigma(t,\cdot,\psi_1,\mu_1)-
\sigma(t,\cdot,\psi_2,\mu_2)
\|^2_{
\ell^2(\mathbb{N},
\mathbb{H})}
\nonumber\\
&=M\sum_{k\in\mathbb{N}}\int_{\mathbb{R}^d}
|\sigma_k(t,x,\psi_1(x),\mu_1)
-\sigma_k(t,x,\psi_2(x),\mu_2)|^2dx
\nonumber\\
&
\leq M\sum_{k\in\mathbb{N}}\int_{\mathbb{R}^d}
\Big(\sigma^2_{4,k}(x)\big(1+ |\psi_1(x)|^{p^*-2}+ |\psi_2(x)|^{p^*-2} +|\psi_1(x)|^{q^*-2}
\nonumber\\& \ \ +|\psi_2(x)|^{q^*-2}
\big)|\psi_1(x)-\psi_2(x)|^2 +\sigma_{5,k}(x)
d_{\mathcal{P}_2}^2(\mu_1,\mu_2)\Big)dx
\nonumber\\
&\leq
\sum_{k\in\mathbb{N}}
\Big(\|\sigma_{4,k}\|^2_{\infty}+
C_3(M,\varepsilon_1)
\|\sigma_{4,k}\|_{\infty}^{\frac{2p-p^*-2}{p-p^*}}
+\widehat{C}_3(M,\varepsilon_2)
\|\sigma_{4,k}\|_{\infty}^{\frac {2q-q^*-2}{q-q^*}}\Big)
\|\psi_1-\psi_2\|^2
\nonumber\\
&\  +M\sum_{k\in\mathbb{N}}
\|\sigma_{5,k}\|_{1}
d_{\mathcal{P}_2}^2(\mu_1,\mu_2)
\nonumber\\
&\  + \frac{\varepsilon_1}{1+ \sum_{k\in\mathbb{N}}
\|\sigma_{4,k}\|_{\infty}}\sum_{k\in\mathbb{N}}
\int_{\mathbb{R}^d}
\sigma_{4,k}(x)\big(|\psi_1(x)|^{p-2}+ |\psi_2(x)|^{p-2}\big)|\psi_1(x)-\psi_2(x)|^2dx
\nonumber\\
&\ + \frac{\varepsilon_2}{1+ \sum_{k\in\mathbb{N}}
\|\sigma_{4,k}\|_{\infty}}
\sum_{k\in\mathbb{N}}\int_{\mathbb{R}^d}
\sigma_{4,k}(x)\big(|\psi_1(x)|^{q-2}+ |\psi_2(x)|^{q-2}\big)|\psi_1(x)-\psi_2(x)|^2dx
\nonumber\\&\leq C_4(M,\varepsilon_1,\varepsilon_2)
\|\psi_1-\psi_2\|^2
 +C_5(M)
d_{\mathcal{P}_2}^2(\mu_1,\mu_2)
\nonumber\\
&\ +\varepsilon_1\int_{\mathbb{R}^d}
 \big(|\psi_1(x)|^{p-2} +|\psi_2(x)|^{p-2}
\big)|\psi_1(x)-\psi_2(x)|^2dx
\nonumber\\
&\  +\varepsilon_2\int_{\mathbb{R}^d}
 \big(|\psi_1(x)|^{q-2} +|\psi_2(x)|^{q-2}
\big)|\psi_1(x)-\psi_2(x)|^2dx,
\end{align}

\begin{align}\label{ffnewA2a0}
\|\mathfrak{L}_{H,S}(t,\psi,\mu)
\|^2_{\mathcal{L}_2(\ell^2
,\mathbb{H})}
 =\|\sigma(t,\cdot,\psi,\mu)
\|^2_{
\ell^2(\mathbb{N},
\mathbb{H})}.
\end{align}

\begin{align}\label{h2----0}
\|\mathfrak{L}_{H,S}(t,\psi_1,\mu_1)-
\mathcal{L}_{H,S}(t,\psi_2,\mu_2)
\|^2_{\mathcal{L}_2(\ell^2
,\mathbb{H})}=\|\sigma(t,\cdot,\psi_1,
\mu_1)-
\sigma(t,\cdot,\psi_2,\mu_2)
\|^2_{
\ell^2(\mathbb{N},
\mathbb{H})}.
\end{align}

\begin{equation}\label{Int1--}
\left\{
  \begin{aligned}
  &d \psi(t)
=\sum_{i=1}^2\textbf{A}_i(t,\psi(t),\mathcal{L}_{
\psi(t)})dt
+ \sqrt{\epsilon}\mathfrak{L}_{H,S}(t,\psi(t),
\mathcal{L}_{
\psi(t)})d\mathcal{W}(t), \ \ t>0, \\
  &\psi(0)=\psi_0\in \mathbb{H},\\
  \end{aligned}
\right.
\end{equation}

\begin{definition}\label{Defsolutionsolution}
(\texttt{Definition of solutions}) Given
$\psi_0\in L^2(\Omega,\mathfrak{F}_0, \mathbb{H})$, a continuous $\mathbb{H}$-valued
$\mathfrak{F}_t$-adapted stochastic process $\psi$
is called a solution to \eqref{Int1--}
if
\begin{align}&\label{Defsolution0}\psi
 \in L^2(\Omega, C([0,T],\mathbb{H} ))\bigcap L^p([0,T]\times\Omega, \mathbb{V}_1)\bigcap L^q([0,T]\times\Omega, \mathbb{V}_2  ),\ \ \forall\ T>0,
\end{align}
and for all $t\geq0$ and $\xi\in\mathbb{H}\cap\mathbb{V}_1
\cap\mathbb{V}_2$, $\mathbf{P}$-a.s.,
\begin{align}\label{Defsolution1}
&\langle \psi(t) ,\xi\rangle=\langle \phi_0 ,\xi\rangle+\sum_{i=1}^2
 \int_{0}^t\ _{\mathbb{V}_i^*}\big\langle \mathbf{A}_i(r,\psi(r),\mathcal{L}_{
\psi(r)}),
\xi\big  \rangle_{\mathbb{V}_i}        dr
 +\int_{0}^t
 \big\langle \xi, \mathfrak{L}_{H,S}(r,\psi(r),\mathcal{L}_{
\psi(r)})d\mathcal{W}(r) \big\rangle.
\end{align}
\end{definition}

\begin{equation}\label{Int1--++}
\left\{
  \begin{aligned}
  &\frac{d \psi_u(t)}{dt}
=\sum_{i=1}^2\textbf{A}_i(t,\psi_u(t),\delta_{\psi^0(t)})
+ \mathfrak{L}_{H,S}(t,\psi_u(t),
\delta_{\psi^0(t)})u(t), \ \ t>0, \\
  &\psi_u(0)=\psi_0\in \mathbb{H},\\
  \end{aligned}
\right.
\end{equation}

\begin{equation}\label{Int1--++++}
\left\{
  \begin{aligned}
  &\frac{d \psi^0(t)}{dt}
=\sum_{i=1}^2\textbf{A}_i(t,\psi^0(t),\delta_{\psi^0(t)})
, \ \ t>0, \\
  &\psi^0(0)=\psi_0\in \mathbb{H}.\\
  \end{aligned}
\right.
\end{equation}

\begin{equation}\label{Int1---}
\left\{
  \begin{aligned}
  &d \psi_u^\epsilon(t)+ \mathfrak{L}_{\mathcal{K}_{p}^\alpha} \psi_u^\epsilon(t) dt
=\sum_{i=1}^2\textbf{A}_i(t,\psi_u^\epsilon(t),
\mathcal{L}_{
\psi^\epsilon(t)})dt\\&+
\mathfrak{L}_{H,S}
(t,\psi_u^\epsilon(t),
\mathcal{L}_{
\psi^\epsilon(t)})u(t)dt+ \sqrt{\epsilon}\mathfrak{L}_{H,S}
(t,\psi_u^\epsilon(t),
\mathcal{L}_{
\psi^\epsilon(t)})d\mathcal{W}(t), \ \ t>0, \\
  &\psi_u^\epsilon(0)=\psi_0\in \mathbb{H}.\\
  \end{aligned}
\right.
\end{equation}

\begin{lemma}\label{SCE1}
Let conditions  {\bf A} and {\bf B1}
be satisfied. Then for every $T>0$, $R>0$ and $N>0$, there exists $C_3(T,R,N)>0$ such that for all $\psi_{0}\in \overline{B}_{R}(\mathbb{H})$, $u\in \mathfrak{A}_N$ and $\epsilon\in(0,1)$,
the solution $\psi^\epsilon_u=\mathfrak{M}_{\psi_0,\mathcal{L}_{\psi^\epsilon}}^\epsilon
\Big(\mathcal{W}+\frac{1}{\sqrt{\epsilon}}
\int_0^\cdot u(t)dt\Big)$ of \eqref{Int1---} satisfies
\begin{align*}
&\mathbf{E}\Big[\|\psi^\epsilon_u
\|^2_{C([0,T],\mathbb{H} )}\Big]+\mathbf{E}\Big[
\|\psi^\epsilon_u\|^p_{L^p([0,T],\mathbb{V}_1 )}\Big]+\mathbf{E}\Big[\|\psi^\epsilon_u\|^q_{L^q([0,T], \mathbb{V}_2)}\Big]\leq C_3(T,R, N),
\end{align*}
\end{lemma}

\begin{lemma}\label{SCE2}
Let conditions  {\bf A} and {\bf B1}
be satisfied. Then for every $T>0$ and $R>0$,
the solution  $\psi^\epsilon$ of \eqref{Int1--} and the solution $\psi^0$
of \eqref{Int1--++++}  satisfy that, for all $\psi_{0}\in \overline{B}_{R}(\mathbb{H})$ and $\epsilon\in(0,1)$,
\begin{align*}
&\mathbf{E}\Big[\|\psi^\epsilon-\psi^0
\|^2_{C([0,T],\mathbb{H} )}\Big]
+\mathbf{E}\Big[
\|\psi^\epsilon-\psi^0\|^p_{L^p([0,T],\mathbb{V}_1 )}\Big]+\mathbf{E}\Big[\|\psi^\epsilon-\psi^0\|^q_{L^q([0,T], \mathbb{V}_2)}\Big]
\leq \epsilon C_4(T,R),
\end{align*}
where $C_4(T,R)>0$ is a constant
independent of $\epsilon$.
\end{lemma}

\begin{lemma}\label{SCE3}
Let conditions  {\bf A} and {\bf B1}
be satisfied. Then for every $T>0$, $\delta>0$, $R>0$ and $N>0$, we have
\begin{align*}
&\lim_{\epsilon\rightarrow0}\sup_{\psi_{0}\in \overline{B}_{R}(\mathbb{H})}
\sup_{u\in \mathfrak{A}_N}\mathbf{P}\Big\{ \|\psi^\epsilon_u-\psi_u\|_{C([0,T],\mathbb{H}) \bigcap L^p([0,T], \mathbb{V}_1)\bigcap L^q([0,T], \mathbb{V}_2)}>\delta\Big\}=0,
\end{align*}
where $\psi^\epsilon_u=
\mathfrak{M}_{\psi_0,\mathcal{L}_{
\psi^\epsilon}}^\epsilon
\big(\mathcal{W}+\frac{1}{\sqrt{\epsilon}}
\int_0^\cdot u(t)dt\big)$ and $\psi_u=
\mathfrak{M}_{\psi_0}^0
\big(
\int_0^\cdot u(t)dt\big)$ are solutions of \eqref{Int1---} and \eqref{Int1--++}, respectively.
\end{lemma}

\begin{proof}
Since $\psi^\epsilon_u=
\mathfrak{M}_{\psi_0,\mathcal{L}_{
\psi^\epsilon}}^\epsilon
\big(\mathcal{W}+\frac{1}{\sqrt{\epsilon}}
\int_0^\cdot u(t)dt\big)$ and $\psi_u=
\mathfrak{M}_{\psi_0}^0
\big(
\int_0^\cdot u(t)dt\big)$ are solutions of \eqref{Int1---} and \eqref{Int1--++}, respectively,
it follows from  It\^{o}'s
formula and
 Lemma \ref{Monotonicity} that for all $t\in[0,T]$, $\psi_{0}\in \overline{B}_{R}(\mathbb{H})$,
$u\in \mathfrak{A}_N$ and $\epsilon\in(0,1)$,
\begin{align}\label{SCE5}
&\|\psi^\epsilon_u(t)
-\psi_u(t)
\|^2+2^{-p}K^{-1}\int_0^{t}
\|\psi^\epsilon_u(r)
-\psi_u(r)\|^p_{\dot{\mathbb{V}}_1}dr
\nonumber\\&\ +2^{1-p}\beta
\int_0^{t}
\|\psi^\epsilon_u(r)
-\psi_u(r)\|_{p}^{p}dr
+2^{1-q}\widehat{\beta}
\int_0^{t}\|\psi^\epsilon_u(r)
-\psi_u(r)\|_{\mathbb{V}_2}^{q}dr
\nonumber\\&\ +\frac{\beta}{2}\int_0^{t} \int_{\mathbb{R}^d}  \big(|\psi^\epsilon_u(r,x)|^{p-2}
+|\psi_u(r,x)|^{p-2}\big)
|\psi^\epsilon_u(r,x)-\psi_u(r,x)|^2dxdr
\nonumber\\&\ +\frac{\widehat{\beta}}{2}
\int_0^{t} \int_{\mathbb{R}^d}  \big(|\psi^\epsilon_u(r,x)|^{q-2}
+|\psi_u(r,x)|^{q-2}\big)
|\psi^\epsilon_u(r,x)-\psi_u(r,x)|^2dxdr
\nonumber\\&\leq
\underbrace{2
\int_0^{t}\big\langle
    \big(\mathfrak{L}_{H,S}(r,\psi^\epsilon_u(r),
\mathcal{L}_{
\psi^\epsilon(t)}) - \mathfrak{L}_{H,S}(r,\psi_u(r),
\delta_{\psi^0(r)})\big)u(r) ,\psi^\epsilon_u(r)
-\psi_u(r)
\big\rangle dr}_{\textbf{L}^\epsilon_9(t)}
\nonumber\\&\ +\underbrace{\epsilon\int_0^{t}
\|\mathfrak{L}_{H,S}(r,\psi^\epsilon_u(r),
\mathcal{L}_{
\psi^\epsilon(r)})
		\|^2_{\mathcal{L}_2(\ell^2,
			\mathbb{H})}dr}_{\textbf{L}^\epsilon_{10}(t)}
\nonumber\\&\ +\underbrace{2\sqrt{\epsilon}
\int_0^{t}\langle
    \psi^\epsilon_u(r)
-\psi_u(r),\mathfrak{L}_{H,S}(r,
\psi^\epsilon_u(r),
\mathcal{L}_{
\psi^\epsilon(r)})d\mathcal{W}(r)
\rangle}_{\textbf{L}^\epsilon_{11}(t)}
\nonumber\\&\ +
2\int_0^{t}
\big(\|\phi_6(r)\|_{\infty}+
\|\widehat{\phi}_6(r)\|_{\infty}\big)
\| \psi^\epsilon_u(r)
-\psi_u(r)
\|^2dr
\nonumber\\&\ +2\int_0^{t}
\big(\|\phi_7(r)\|_{1}+
\|\widehat{\phi}_7(r)\|_{1}\big)
d_{\mathcal{P}_2}^2
(\mathcal{L}_{
\psi^\epsilon(r)},\delta_{\psi^0(r)})dr.
\end{align}
For   $\textbf{L}^\epsilon_{9}(t)$ in \eqref{SCE5}, by \eqref{h2--} (for $M=1$, $\varepsilon_1=\beta/2$ and
$\varepsilon_2=\widehat{\beta}/2$) and \eqref{h2----0},
we deduce that for all $t\in[0,T]$, $\psi_{0}\in \overline{B}_{R}(\mathbb{H})$,
$u\in \mathfrak{A}_N$ and $\epsilon\in(0,1)$,
\begin{align}
\label{SCE6}
\textbf{L}^\epsilon_{9}(t)
&\leq
\int_0^{t}
\|u(r)\|^2\|\psi^\epsilon_u(r)
-\psi_u(r)
\|^2dr\nonumber\\&\ +
\int_0^{t}
    \big\|\mathfrak{L}_{H,S}(r,\psi^\epsilon_u(r),
\mathcal{L}_{
\psi^\epsilon(t)}) - \mathfrak{L}_{H,S}(r,\psi_u(r),
\delta_{\psi^0(r)})
\big\|^2 dr
\nonumber\\&\leq K_1
\int_0^{t}
\big(1
+\|u(r)\|^2\big)\|\psi^\epsilon_u(r)
-\psi_u(r)
\|^2dr
+ K_1\int_0^{t}
d_{\mathcal{P}_2}^2
(\mathcal{L}_{
\psi^\epsilon(r)},\delta_{\psi^0(r)})dr
\nonumber\\&\ +\frac{\beta}{2}\int_0^{t} \int_{\mathbb{R}^d}  \big(|\psi^\epsilon_u(r,x)|^{p-2}
+|\psi_u(r,x)|^{p-2}\big)
|\psi^\epsilon_u(r,x)-\psi_u(r,x)|^2dxdr
\nonumber\\&\ +\frac{\widehat{\beta}}{2}
\int_0^{t} \int_{\mathbb{R}^d}  \big(|\psi^\epsilon_u(r,x)|^{q-2}
+|\psi_u(r,x)|^{q-2}\big)
|\psi^\epsilon_u(r,x)-\psi_u(r,x)|^2dxdr,
\end{align}
where $K_1>0$ is a constant independent of $\epsilon$.

By \eqref{SCE5} and \eqref{SCE6}, we have
\begin{align}\label{SCE5-1}
&\sup_{r\in [0,t]}\|\psi^\epsilon_u(r)-\psi_u(r)\|^2
+2^{-p}K^{-1}\int_0^{t}
\|\psi^\epsilon_u(r)
-\psi_u(r)\|^p_{\dot{\mathbb{V}}_1}dr
\nonumber\\&\ +2^{1-p}\beta
\int_0^{t}
\|\psi^\epsilon_u(r)
-\psi_u(r)\|_{p}^{p}dr
+2^{1-q}\widehat{\beta}
\int_0^{t}\|\psi^\epsilon_u(r)
-\psi_u(r)\|_{\mathbb{V}_2}^{q}dr \nonumber\\
&\leq
2\int_0^{t}
\big(K_1
+K_1\|u(r)\|^2
+2\|\phi_6(r)\|_{\infty}+
2\|\widehat{\phi}_6(r)\|_{\infty}\big)\|\psi^\epsilon_u(r)
-\psi_u(r)
\|^2dr
\nonumber\\&\ + 2\int_0^{t}
\big(K_1+2 \|\phi_7(r)\|_{1}+
2\|\widehat{\phi}_7(r)\|_{1}\big)
d_{\mathcal{P}_2}^2
(\mathcal{L}_{
\psi^\epsilon(r)},\delta_{\psi^0(r)})dr
+2\textbf{L}^\epsilon_{10}(t)
+2\sup_{r\in[0,t]}|\textbf{L}^\epsilon_{11}(r)|.
\end{align}
Then
\begin{align}\label{SCE5-2}
&\sup_{r\in [0,t]}\|\psi^\epsilon_u(r)-\psi_u(r)\|^2
+2^{-p}K^{-1}\int_0^{t}
\|\psi^\epsilon_u(r)
-\psi_u(r)\|^p_{\dot{\mathbb{V}}_1}dr
\nonumber\\&\ +2^{1-p}\beta
\int_0^{t}
\|\psi^\epsilon_u(r)
-\psi_u(r)\|_{p}^{p}dr
+2^{1-q}\widehat{\beta}
\int_0^{t}\|\psi^\epsilon_u(r)
-\psi_u(r)\|_{\mathbb{V}_2}^{q}dr \nonumber\\
&\leq
e^{2\int_0^{t}
\big(K_1
+K_1\|u(r)\|^2
+2\|\phi_6(r)\|_{\infty}+
2\|\widehat{\phi}_6(r)\|_{\infty}\big)dr}
\nonumber\\&\ \ \cdot\Big( 2\int_0^{t}
\big(K_1+2 \|\phi_7(r)\|_{1}+
2\|\widehat{\phi}_7(r)\|_{1}\big)
d_{\mathcal{P}_2}^2
(\mathcal{L}_{
\psi^\epsilon(r)},\delta_{\psi^0(r)})dr
+2\textbf{L}^\epsilon_{10}(t)
+2\sup_{r\in[0,t]}|\textbf{L}^\epsilon_{11}(r)|
\Big)
\nonumber\\
&\leq
e^{2
\big(K_1 T
+K_1 N^2
+2 \int_0^T \left(\|\phi_6(r)\|_{\infty}+
\|\widehat{\phi}_6(r)\|_{\infty}\right)dr \big)}
\nonumber\\&\ \ \cdot \Big(2\int_0^{t}
\big(K_1+2 \|\phi_7(r)\|_{1}+
2\|\widehat{\phi}_7(r)\|_{1}\big)dr
\mathbf{E}\Big[
\| \psi^\epsilon-\psi^0\|^2_{C([0,T],\mathbb{H} )}\Big]
+2\textbf{L}^\epsilon_{10}(t)
+2\sup_{r\in[0,t]}|\textbf{L}^\epsilon_{11}(r)|
\Big).
\end{align}

For the  terms $\textbf{L}^\epsilon_{10}(t)$ and $\sup_{r\in[0,t]}|\textbf{L}^\epsilon_{11}(r)|$ in \eqref{SCE5-2}, by \eqref{h2-}, \eqref{ffnewA2a0}
and the BDG inequality,
we infer that for all $t\in[0,T]$, $\psi_{0}\in \overline{B}_{R}(\mathbb{H})$,
$u\in \mathfrak{A}_N$ and $\epsilon\in(0,1)$,
\begin{align}&
\label{SCE7}
2 e^{2
\big(K_1 T
+K_1 N^2
+2 \int_0^T \left(\|\phi_6(r)\|_{\infty}+
\|\widehat{\phi}_6(r)\|_{\infty}\right)dr \big)}
\left( \mathbf{E}[
\textbf{L}^\epsilon_{10}(t)]+
\mathbf{E}\Big[\sup_{r\in[0,t]}
|\textbf{L}^\epsilon_{11}(r)|\Big]
\right)
\nonumber\\&\leq\frac{1}{2}
\mathbf{E}\Big[\sup_{r\in[0,t]}
\|\psi^\epsilon_u(r)
-\psi_u(r)
\|^2\Big]+\epsilon K_2\mathbf{E}\bigg[\int_0^{t}
\|\mathfrak{L}_{H,S}(r,\psi^\epsilon_u(r),
\mathcal{L}_{
\psi^\epsilon(r)})
		\|^2_{\mathcal{L}_2(\ell^2,
			\mathbb{H})}dr\bigg]
\nonumber\\&\leq \frac{1}{2}\mathbf{E}\Big[\sup_{r\in[0,t]}
\|\psi^\epsilon_u(r)
-\psi_u(r)
\|^2\Big]+\epsilon K_3
\int_0^{t}
\Big(1+\textbf{E}[\|\psi^\epsilon_u(r)\|_{p
}^{p}]+\textbf{E}[\|\psi^\epsilon_u(r)
\|_{\mathbb{V}_2}^{q}]+
\textbf{E}[\|\psi^\epsilon(r)\|^2]\Big)dr,
\end{align}where $K_2$ and $K_3$ are constants independent of $\epsilon$.
It follows from \eqref{SCE5-2} and \eqref{SCE7}
 that for all $t\in[0,T]$, $\psi_{0}\in \overline{B}_{R}(\mathbb{H})$,
$u\in \mathfrak{A}_N$ and $\epsilon\in(0,1)$,
 \begin{align}\label{SCE8}
&\mathbf{E}\Big[\sup_{r\in[0,t]}
\|\psi^\epsilon_u(r)
-\psi_u(r)
\|^2\Big]+2^{1-p}K^{-1}\mathbf{E}
\bigg[\int_0^{t}
\|\psi^\epsilon_u(r)
-\psi_u(r)\|^p_{\dot{\mathbb{V}}_1}dr\bigg]
\nonumber\\&\ +2^{2-p}\beta\mathbf{E}
\bigg[
\int_0^{t}
\|\psi^\epsilon_u(r)
-\psi_u(r)\|_{p}^{p}dr\bigg]
+2^{2-q}\widehat{\beta}\mathbf{E}
\bigg[
\int_0^{t}\|\psi^\epsilon_u(r)
-\psi_u(r)\|_{\mathbb{V}_2}^{q}dr\bigg]
\nonumber\\&\leq
K_4 \int_0^{t}
\big(K_1+2 \|\phi_7(r)\|_{1}+
2\|\widehat{\phi}_7(r)\|_{1}\big)dr
\mathbf{E}\Big[
\| \psi^\epsilon-\psi^0\|^2_{C([0,T],\mathbb{H} )}\Big]
\nonumber\\&\ +\epsilon K_4
\int_0^{t}
\Big(1+\textbf{E}[\|\psi^\epsilon_u(r)\|_{p
}^{p}]+\textbf{E}[\|\psi^\epsilon_u(r)
\|_{\mathbb{V}_2}^{q}]+
\textbf{E}[\|\psi^\epsilon(r)\|^2]\Big)dr,
\end{align}
where $K_4>0$ is a constant independent of $\epsilon$. Then we find from \eqref{SCE8}, Theorem \ref{Theorem4} and Lemma \ref{SCE1} that
for all $t\in [0,T]$,
 \begin{align}\label{SCE9}
&\mathbf{E}\Big[\sup_{r\in[0,t]}
\|\psi^\epsilon_u(r)
-\psi_u(r)
\|^2\Big]+2^{1-p}K^{-1}\mathbf{E}
\bigg[\int_0^{t}
\|\psi^\epsilon_u(r)
-\psi_u(r)\|^p_{\dot{\mathbb{V}}_1}dr\bigg]
\nonumber\\&\ +2^{2-p}\beta\mathbf{E}
\bigg[
\int_0^{t}
\|\psi^\epsilon_u(r)
-\psi_u(r)\|_{p}^{p}dr\bigg]
+2^{2-q}\widehat{\beta}\mathbf{E}
\bigg[
\int_0^{t}\|\psi^\epsilon_u(r)
-\psi_u(r)\|_{\mathbb{V}_2}^{q}dr\bigg]
\nonumber\\&\leq
K_5\mathbf{E}\Big[\|\psi^\epsilon-\psi^0
\|^2_{C([0,T],\mathbb{H} )}\Big]
+\epsilon K_5,
\end{align}
where $K_5=K_5(T,N,R)>0$  is a constant independent of $\epsilon$.

By \eqref{SCE9} and Lemma \ref{SCE2}, we obtain
\begin{align}\label{SCE10}
&\lim_{\epsilon\rightarrow0}\sup_{\psi_{0}\in \overline{B}_{R}(\mathbb{H})}
\sup_{u\in \mathfrak{A}_N}
\mathbf{E}\Big[\sup_{r\in[0,T]}
\|\psi^\epsilon_u(r)
-\psi_u(r)
\|^2\Big]=0,
\\&\label{SCE11}\lim_{\epsilon\rightarrow0}
\sup_{\psi_{0}\in \overline{B}_{R}(\mathbb{H})}
\sup_{u\in \mathfrak{A}_N}\mathbf{E}
\bigg[\int_0^{T}
\|\psi^\epsilon_u(r)
-\psi_u(r)\|^p_{\mathbb{V}_1}dr\bigg]
=0,
\end{align}
and
\begin{align}\label{SCE12}
\lim_{\epsilon\rightarrow0}\sup_{\psi_{0}\in \overline{B}_{R}(\mathbb{H})}
\sup_{u\in \mathfrak{A}_N}\mathbf{E}
\bigg[
\int_0^{T}\|\psi^\epsilon_u(r)
-\psi_u(r)\|_{\mathbb{V}_2}^{q}dr\bigg]=0.
\end{align}
As a result of \eqref{SCE10}-\eqref{SCE12}, we find that
$$\lim_{\epsilon\rightarrow0}\sup_{\psi_{0}\in \overline{B}_{R}(\mathbb{H})}
\sup_{u\in \mathfrak{A}_N}
\mathbf{E}\Big[ \|\psi^\epsilon_u-\psi_u\|^2_{C([0,T],\mathbb{H}) \bigcap L^p([0,T], \mathbb{V}_1)\bigcap L^q([0,T], \mathbb{V}_2)}\Big]=0.$$
From this and Chebyshev inequality, we conclude the proof.
\end{proof}

\end{document}
